\documentclass[12pt,addpoints]{evalua}
\grado{2$^\circ$ de Secundaria}
\cicloescolar{2025-2026}
\materia{Matemáticas 2}
\unidad{1}
\title{Examen Extraordinario}
\aprendizajes{
      \item Resuelve problemas de multiplicación y división con números enteros, fracciones y decimales positivos y negativos.
      \item Resuelve problemas de potencias con exponente entero y aproxima raíces cuadradas.
      \item Resuelve problemas que impliquen el uso de la notación científica.
      \item Calcula porcentajes de cantidades.
}
\author{Prof.: Julio César Melchor Pinto}
\begin{document}
\begin{multicols}{2}
	\tableofcontents
\end{multicols}\newpage
\begin{questions}
      \section{Cálculos numéricos}
           \question[3]{Realiza las siguientes operaciones de \textit{cálculo numérico}:

           \begin{parts}
                       \begin{multicols}{2}
            	\subsection{Suma de números}
                  \part \include*{../questions/ques_suma_numeros_1}
	\subsection{Resta de números}
                  \part \include*{../questions/ques_resta_numeros_1}
	\subsection{Multiplicación de números}
                  \part \include*{../questions/ques_multiplicacion_numeros_1}
	\subsection{División de números}
                  \part \include*{../questions/ques_division_numeros_1}
      \end{multicols}
      \subsection{Resolución de Problemas}
                   \part \include*{../questions/ques_resolucion _problemas_1}
\end{parts}
}


\section{Fracciones}

\subsection{Clasificación de fracciones}
      \question[3]{}
\subsection{Representación de fracciones}
      \question[3]{}
\subsection{Nombre de fracciones}
      \question[3]{}
\subsection{Fracciones en la recta numérica}
      \question[3]{}
\subsection{Conversión de fracciones}
      \question[3]{}

\section{Fracciones MCD y MCM}
 \subsection{Comparación de fracciones}
      \question[3]{}

\subsection{Fracciones equivalentes}
      \question[3]{}

\subsection{M.C.D y M.C.M}
      \question[3]{}

\subsection{Simplificación de fracciones}
      \question[3]{}

\subsection{Resolución de problemas}
      \question[3]{}

\section{Números decimales}
      \subsection{Ubicación en la recta numérica}
      \question[3]{}
\subsection{Porcentajes a decimal}
      \question[3]{}
\subsection{Operaciones con múltiplos de 10}
      \question[3]{}
\subsection{Conversión de fracciones a decimales}
      \question[3]{}
\subsection{Conversión de decimales a fracciones}
      \question[3]{}

\section{Conceptos básicos de números negativos}
\subsection{Ubicación en la recta numérica}
      \question[3]{}
\subsection{Comparación de negativos}
      \question[3]{}
\subsection{Oraciones con negativos}
      \question[3]{}
\subsection{Determina el signo en suma y resta con negativos}
      \question[3]{}
\subsection{Determina el signo multiplicación y división con negativos}
      \question[3]{}

\section{Operaciones con números negativos}
\subsection{Comparación de negativos}
      \question[3]{}
\subsection{Suma y resta con negativos}
      \question[3]{}
\subsection{Multiplicación y división con negativos}
      \question[3]{}
\subsection{Potencias con numeros negativos}
      \question[3]{}
\subsection{Jerarquía de operacione}
      \question[3]{}

\end{questions}
\end{document}